\section{Every characteristic but three}\label{sec:fields}

A code is described by finite data over a commutative ring, and the data are moved to a field by a ring homomorphism.

\begin{proposition}[Certificates]\label{prop:certificate}
\formal{Transfer.lean}{Certificate,Certificate.RowsHold,Certificate.PivotsHold,Certificate.code,Certificate.dec,Certificate.row_eq_dec_of_mem,Certificate.isClassUniformCode,Certificate.finrank_code}
Let $R$ be a commutative ring and $\varphi\colon R\to F$ a ring homomorphism to a field. Fix a colouring of $\Z/d$,
matrices $G_1,\dots,G_m\in\Mat_d(R)$, for every class $J$ and every $r\in\Z/d$ coefficients $a_{J,r}(i,j)\in R$,
elements $N,P\in R$, and positions $p_1,\dots,p_m$, such that
\begin{enumerate}[label=(\roman*),leftmargin=2em]
\item $N\,(G_k)_{c,r}=\sum_{i,j}a_{J,r}(i,j)\,(G_k)_{i,j+c}$ for every $k$, every $r$, and every $c$ in a class $J$;
\item $(G_k)_{p_l}=P$ if $k=l$ and $(G_k)_{p_l}=0$ otherwise.
\end{enumerate}
If $\varphi(N)\ne0$, the span of the matrices $\varphi(G_k)$ is a code of the class-uniform rule
$\cA_J(Y)_r=\varphi(N)^{-1}\sum_{i,j}\varphi(a_{J,r}(i,j))\,Y_{i,j}$. If moreover $\varphi(P)\ne0$, it has
dimension~$m$.
\end{proposition}

\begin{outline}
Applying $\varphi$ to~(i) and dividing by $\varphi(N)$ shows that each $\varphi(G_k)$ satisfies the rule
(\lean{Certificate.genMap_row}); the rule is linear, so the span satisfies it (\lean{Certificate.row_eq_dec_of_mem}).
The map $Y\mapsto(\varphi(P)^{-1}Y_{p_l})_l$ sends $\varphi(G_k)$ to the $k$-th standard basis vector, so the
$\varphi(G_k)$ are linearly independent (\lean{Certificate.linearIndependent_genMap}).
\end{outline}

\begin{theorem}[Theorem~A, every characteristic but three]\label{thm:fields}
\formal{Fields.lean}{exists_code_of_three_ne_zero,exists_classUniformCode_gt_of_three_ne_zero,exists_code_of_charZero,exists_code_zmod}
Let $F$ be a field in which $3\ne0$. Then the classes $\{0,1,3\}$ and $\{2,4,5\}$ of $\Z/6$ carry a code of a
class-uniform rule of dimension $25$ over $F$, so that $2\dim\X=50>48=6^2+2\cdot6$. In particular this holds over
every field of characteristic zero and over $\F_p$ for every prime $p\ne3$.
\end{theorem}

\begin{outline}
Three certificates in the sense of \cref{prop:certificate}, all for the parent paper's colouring
\lean{SwapRigidity.aperiodicCls}, cover the fields.
\begin{itemize}[leftmargin=1.5em]
\item \lean{rationalA}, over $R=\Z$ with $\varphi$ the canonical map, has $N=1008=2^4\cdot3^2\cdot7$ and $P=3$. It serves
every field in which $2$, $3$ and $7$ are nonzero.
\item \lean{rationalB}, over $\Z$, has $N=187200=2^6\cdot3^2\cdot5^2\cdot13$ and $P=15$. In characteristic $7$ its scales
are $187200=6$ and $15=1$, both nonzero; it serves characteristic~$7$.
\item \lean{binary}, over $R=\F_2$, is the parent paper's code with $N=P=1$ (\lean{binary_rowsHold},
\lean{binary_pivotsHold}); with $\varphi$ the map $\F_2\to F$ it serves characteristic~$2$.
\end{itemize}
The conditions (i) and~(ii) for the two rational certificates, $900$ identities between integers and $625$ entries each,
are proved by \texttt{decide} (\lean{rationalA_rowsHold}, \lean{rationalA_pivotsHold}, \lean{rationalB_rowsHold},
\lean{rationalB_pivotsHold}). The proof of \lean{exists_code_of_three_ne_zero} splits on whether $2=0$ in $F$, then on
whether $7=0$; in the remaining case $1008=2^4\cdot3^2\cdot7\ne0$. For $\F_p$, $3=0$ only when $p$ divides~$3$.
\end{outline}

\begin{remark}[Numerical: the certificates]\label{rem:certificates}
The data are in \leanfile{tools/data/rational_a.json} and \leanfile{tools/data/rational_b.json}, each a class-uniform
rule over $\Q$. The generator \leanfile{tools/gen_rational_data.py} brings the $36\times36$ system of the rule to
reduced echelon form, takes the standard basis of its solution space (equal to $1$ at one free position and $0$ at the
others), clears the denominators of the basis and of the decoders, and writes the integer data to
\leanfile{RationalData.lean}. Lean does not trust the generator: the kernel checks the identities (i) and (ii), and the
proofs use nothing else. Two further facts are checked by exact computation only, by the generator and by an
independent recomputation, and no proof uses them: the solution space of each rule has dimension exactly $25$, so the
span of the generators is the whole code of the rule, and the two codes are different subspaces of $\Mat_6(\Q)$,
meeting in dimension $20$. Between them the two rational certificates serve every characteristic except $2$ and~$3$.
\end{remark}

The parent paper records that on cone relations, linear recovery with bounded complex coefficients and bounded width
stays at the pairing rate, so that a linear construction needs positive characteristic or coefficients that cannot be
rescaled to be bounded \cite[Section~4.5]{SeventeenthPaper2026}. A reader could take the gain of \cref{prop:F2}, the one
gain over the interval value that the parent paper proves, for a phenomenon of characteristic~$2$ of the same kind.
\Cref{thm:fields} shows that it is not. It concerns a single size and says nothing about rates, so it does not bear on
that statement; it removes the only reason the examples gave to restrict a search for larger gains to positive
characteristic. Characteristic~$3$ is open (\cref{prob:char3}).
